\section{Learning to Use the Optional Reference Bridges}
Twenty-one native fits regress two conservative gains and the
components of all/easy harm from observed motion and candidate paths.
All fitting excludes the outer location. Query decisions choose at
most two edits and keep at least one CV fallback. The joint rule uses
predicted constraints; the original recording-level test evaluates
realized forecasts separately.

The joint rule makes \CausalBridgeEdits{} edits versus
\CausalBridgeIndependent{} for independent utility. Its matched-count
conservative difference is \CausalBridgeDirect{} percentage points
and \CausalBridgeUndefined{} views have undefined easy support.
Only \CausalBridgePass/\CausalBridgeViews{} views pass risk and improve
both references. The exact hindsight program at the same query cap
verifies \CausalBridgeOracle/\CausalBridgeViews{}. With the same
learned rule, substituting outcome-derived gain alone gives
\CausalBridgeGain/\CausalBridgeViews{}, while substituting
outcome-derived risk alone gives \CausalBridgeRisk/\CausalBridgeViews{}.
This identifies the predicted risk screen as the dominant bottleneck
in this rule. The substitutions are diagnostic, not causal policies.
